\subsection{Propagation Calculation}
\label{subsec:propagation-calculation}

On a cache miss, ns-3 synchronizes the positions of all UEs to the Python
scene and adds eligible virtual receivers before issuing a single batched ray
tracing call.
Let $\mathcal{R}$ be the physical receiver set and $\mathcal{V}_r$ the virtual
receivers generated for UE $r$; ray tracing is performed over the augmented set
\begin{equation}
    \mathcal{R}^{+}
    =
    \mathcal{R}\cup\bigcup_{r\in\mathcal{R}}\mathcal{V}_r .
\end{equation}
The Sionna \texttt{PathSolver} is called with maximum depth $3$, line-of-sight
paths, specular reflections, refraction, and synthetic arrays enabled; diffuse
reflection is disabled \cite{sionnart,alsamman_mmwave,raymobtime_sionna}.
The result is a single geometrically consistent snapshot of propagation
parameters for all present and predicted receiver positions.

For each Tx–Rx pair $(b,r)$, the solver returns a set of propagation paths
$\{(\tau_{b,r,p}, \mathbf{h}_{b,r,p})\}$ indexed by $p$.
The reported propagation delay is the minimum non-negative path delay,
\begin{equation}
    d_{b,r}=
    \mathrm{round}
    \left(
    10^9 \min_{p:\tau_{b,r,p}\geq 0}\tau_{b,r,p}
    \right)\;\text{[ns]}.
\end{equation}
Combining all paths across the OFDM bandwidth yields the MIMO channel frequency
response
\begin{equation}
    \mathbf{H}_{b,r}\in
    \mathbb{C}^{N_{\mathrm{rx}}\times N_{\mathrm{tx}}\times N_f},
\end{equation}
where $N_{\mathrm{rx}}$, $N_{\mathrm{tx}}$, and $N_f$ are the numbers of
receive antenna elements, transmit antenna elements, and resource blocks,
respectively.

The average channel power and path loss are then computed as
\begin{align}
    P_{b,r} & =
    \frac{1}{N_{\mathrm{rx}}N_{\mathrm{tx}}N_f}
    \sum_{i,j,q}\left|H_{b,r}[i,j,q]\right|^2, \\
    L_{b,r} & =-10\log_{10}(P_{b,r}).
\end{align}
If $P_{b,r}=0$, the implementation exports $L_{b,r}=200$ dB to indicate no signal.
The channel is then normalized as
\begin{equation}
    \widetilde{H}_{b,r}[i,j,q]
    =
    \frac{H_{b,r}[i,j,q]}{\sqrt{P_{b,r}}},
\end{equation}
separating large-scale attenuation from small-scale spatial and frequency
selectivity: $L_{b,r}$ is consumed by the ns-3 \texttt{PropagationLossModel},
while $\widetilde{\mathbf{H}}_{b,r}$ is passed to the spectrum channel matrix.
For ns-3 export, the tensor is serialized according to the \texttt{MatrixArray}
page-major, column-major layout: the flat index is
\begin{equation}
    k = q\,N_{\mathrm{tx}}N_{\mathrm{rx}}
    + j\,N_{\mathrm{rx}}
    + i,
\end{equation}
where $q$, $j$, and $i$ index resource block, transmit antenna element, and
receive antenna element, respectively.

After the records are returned to ns-3, the temporary virtual receivers are
removed from the Python scene.
The ns-3 propagation cache stores all returned records for both real and
virtual receiver positions using the same displacement-validity rule, so that
future queries for nearby positions are served from the cache without invoking
the ray tracer again.
